\documentclass[12pt]{article}
\usepackage[utf8]{inputenc}
\usepackage[T1]{fontenc}
\usepackage{hyperref}
\title{Teiwa Language Resources for K-12 Curriculum Development}
\author{Compiled by Claude Sonnet 5 (claude-sonnet-5) for Andr\'as Kornai}
\date{}
\begin{document}
\maketitle

\section{Introduction}
Teiwa (ISO 639-3 twe; Glottocode teiw1235) is a Timor-Alor-Pantar language spoken on Pantar island in eastern Indonesia. The task description gives about 1,500 speakers; a search summary of the Klamer grammar gave about 4,000. I could not settle the difference. The language is described as highly endangered. The main documentation is academic. I found no school materials.

\section{Primary Educational Resources}
None found. No primers, textbooks or school programmes for Teiwa turned up in search.

\section{Linguistic and Reference Materials}
\begin{itemize}
\item Marian Klamer, \emph{A Grammar of Teiwa} (De Gruyter Mouton, 2010, Mouton Grammar Library 49). It covers phonology, word classes, noun phrases, serial verbs, clauses and information structure. It includes two narrative texts, two songs and an English-Teiwa and Teiwa-English word list of 1,312 items. Publisher page: \url{https://www.degruyterbrill.com/document/doi/10.1515/9783110226072/html}
\item Library record of the same grammar: \url{https://www.oeaw.ac.at/resources/Record/9783110226072}
\item Perlego listing of the grammar: \url{https://www.perlego.com/book/403862/a-grammar-of-teiwa-pdf}
\item LINGUIST List review of the grammar: \url{https://linguistlist.org/issues/22/3994/}
\item Leiden University repository item (not opened): \url{https://scholarlypublications.universiteitleiden.nl/access/item%3A3719511/view}
\end{itemize}

\section{Religious and Community Resources}
None found.

\section{Digital and Audio Resources}
\begin{itemize}
\item Endangered Languages Project entry for Teiwa: \url{https://www.endangeredlanguages.com/lang/2087?hl=pt-br}
\end{itemize}
No audio collections were seen in the search results.

\section{Curriculum Development Context}
Teiwa has one strong reference grammar with a short word list and some texts and songs. Nothing found is aimed at schools. A curriculum would need new material written with Teiwa speakers, using Klamer's word list and texts as a starting point. Contact with the author and the Pantar community would be needed, and I have no contact details to give.
\end{document}